\subsection{An explicit ordinary class character and its balanced image}

\begin{statement}{A.170}{Proved}
{A compatible unramified cubic character with exact source and target}
Retain the tower \(K_n\) of A.165--A.168, put
\[
 F=\Q(\alpha,\zeta),\qquad E_n=K_n(\zeta),\qquad
 \alpha^3=3,\quad \zeta=\zeta_3,\quad d=649,
\]
and let \(\tau\) fix \(K_n\) and invert \(\zeta\).
Set
\[
 \begin{gathered}
 g_2=\alpha-1,\quad g_{11}=2+\alpha,\quad
 g_{59}=2+2\alpha-\alpha^2,\qquad
 \sigma(\alpha)=\zeta\alpha,\quad\sigma(\zeta)=\zeta,\\
 u_q=\frac{\sigma(g_q)}{\sigma^2(g_q)}
       \quad(q=2,11,59),\qquad
 u=u_{11}u_{59}
      =\frac{1+6\zeta\alpha}{1+6\zeta^2\alpha}.
 \end{gathered}
\tag{A.170.1}
\]
Choose \(r^3=u\).  The involution
\(\widetilde\tau(r)=r^{-1}\), extending \(\tau\), defines actual
cyclic cubic extensions
\[
 \mathscr H_n=E_n(r)^{\langle\widetilde\tau\rangle}/K_n
                   \qquad(n\ge2).
\tag{A.170.2}
\]
They are nontrivial, unramified at every finite prime, and satisfy
\(\mathscr H_m=K_m\mathscr H_n\) for \(m\ge n\).
They split at every prime above two and three.
The cyclic cubic polynomial giving this field before base change is
\[
 Z^3-3Z-
     \frac{2-12\alpha-36\alpha^2}{1-6\alpha+36\alpha^2}.
\tag{A.170.3}
\]
Identify its Galois group with \(\Z/3\) by
\(r\mapsto\zeta r\).  Arithmetic Frobenius gives compatible
surjective class characters
\(\chi_n:\operatorname{Cl}(K_n)\{3\}\to\Z/3\).
On the actual changed-prime branches their values are
\[
 \begin{array}{c|rrrr}
 P&P_{11,0,n}&P_{11,2,n}&P_{59,0,n}&P_{59,2,n}\\ \hline
 \chi_n([P])&2&0&1&0 .
 \end{array}
\tag{A.170.4}
\]
In particular \(\chi_n([P_{2,0,n}])=0\), and \(\chi_n\) does not
factor through the balanced quotient \(C_n^{\rm bal}\) of A.155.
Together with A.168 this proves the unconditional size bound
\[
                 |\operatorname{Cl}(K_n)\{3\}|\ge9
                       \qquad(n\ge3).
\tag{A.170.5}
\]
The same bound holds for the subgroup generated by the three-primary
classes of primes over \(S=\{2,3,11,59,\infty\}\).

There is also an exact finite-source classification.  Let
\[
 \mathcal A^-=
  \bigl(\mathcal O_{F,S_F}^{\times}/
                    (\mathcal O_{F,S_F}^{\times})^3\bigr)^{\tau=-1}.
\]
For every \(n\ge2\), the classes in \(\mathcal A^-\) whose actual
Kummer extension over \(E_n\) is unramified at every finite place
form precisely the line
\[
                              \F_3[u].
\tag{A.170.6}
\]
This describes the indicated radical source over \(F\); it does not
exhaust the class characters of \(K_n\).
The balanced class \(e_3\) remains undecided by these results.
\end{statement}

\begin{proof}
\emph{The actual field maps and the first unit.}
The identity
\((2+x)(2+2x-x^2)=1+6x\) for \(x^3=3\) proves the last
equality of A.170.1.  Every \(u_q\) is taken to its inverse by
\(\tau\).  In a nontrivial cubic Kummer extension the specified
\(\widetilde\tau\) has order two and commutes with
\(\kappa:r\mapsto\zeta r\):
\(\widetilde\tau\kappa\widetilde\tau(r)=\zeta r\).
Its fixed field is therefore cyclic cubic over the fixed field of
\(\tau\).  The element \(Z=r+r^{-1}\) lies in that fixed field and
satisfies A.170.3.  The displayed coefficient follows by multiplying
the numerator and denominator in A.170.1:
\[
 \begin{gathered}
 D=1-6\alpha+36\alpha^2,\qquad
 u+u^{-1}=(2-12\alpha-36\alpha^2)/D,\\
 \operatorname{disc}(\text{A.170.3})
       =\left(\frac{108\alpha(1-3\alpha)}{D}\right)^2 .
 \end{gathered}
\tag{A.170.7}
\]
The fixed-field construction supplies the field map; the square
discriminant is an independent exact check.

The following explicit unit will also prove the completeness of
the source in A.170.6:
\[
 \begin{gathered}
 \lambda=\zeta-1,\qquad \pi=\frac{\alpha^2}{\lambda},\qquad
 \pi^6=-3,\quad \pi^3=\zeta^2-\zeta,\quad
 \tau(\pi)=-\zeta\pi,\\
 v=\frac{1+\pi}{1+\pi^2},\qquad
 v\tau(v)=\alpha^2-2=U,\qquad
 \epsilon=\frac{v}{\tau(v)}.
 \end{gathered}
\tag{A.170.8}
\]
All identities follow from \(\alpha^3=3\) and
\(\lambda^2=-3\zeta\).
The inverse expressions
\(\zeta=-(1+\pi^3)/2\) and \(\alpha=-\zeta\pi^2\)
show that \(F=\Q(\pi)\), so \(X^6+3\) is its minimal polynomial.
Both \(1+\pi\) and \(1+\pi^2\) have absolute norm four.
At two the splitting field \(F\) is unramified, with three primes
of residue degree two.  Fixing a residue \(\zeta\in\F_4\), the
three residues \(\alpha=1,\zeta,\zeta^2\) give respectively
\(\pi=\zeta,1,\zeta^2\).
Thus \(1+\pi\) and \(1+\pi^2\) vanish at one and the same prime
of norm four.  Their equal norms give valuation one there and zero at
every other prime.  Their quotient \(v\) is an actual global unit.
The identity for \(v\tau(v)\) in A.170.8 proves in particular that
\(\epsilon\) is a norm-one unit.

\emph{The exact local unit classes at three.}
Write \(F_{\mathfrak p}\) for the completion of \(F\) at its unique
three-prime, and use its normalized valuation
\[
 v_{\mathfrak p}(3)=6,\quad v_{\mathfrak p}(\alpha)=2,\quad
 v_{\mathfrak p}(\lambda)=3,\quad v_{\mathfrak p}(\pi)=1.
\tag{A.170.9}
\]
The residue field is \(\F_3\).  Direct subtraction gives
\[
 v_{\mathfrak p}(\epsilon-1)=1,\qquad
 v_{\mathfrak p}(u_{11}-1)=v_{\mathfrak p}(u_{59}-1)=5,\qquad
                         v_{\mathfrak p}(u-1)=11.
\tag{A.170.10}
\]
For the first equality, \(v\equiv1+\pi\) and
\(\tau(v)\equiv1-\pi\pmod{\pi^2}\).
For the middle equalities the numerators of the differences are
\(\alpha(\zeta-\zeta^2)\) and
\((2\alpha+\alpha^2)(\zeta-\zeta^2)\), with unit denominators.
For the last equality use \(6\alpha(\zeta-\zeta^2)\) in A.170.1.

The last unit is a cube in the local field \(F_{\mathfrak p}\).
Indeed write \(u=1+C\), put \(c=C/3\), and note \(v_{\mathfrak p}(c)=5\).
For \(X^3-u\) at \(X=1+c\), the residual
\(3c^2+c^3\) has valuation fifteen, whereas the derivative has
valuation six.  The strict Hensel inequality \(15>2\cdot6\)
gives the root.

An exact representative of the unramified cubic Kummer line over
\(F_{\mathfrak p}\) is
\(1+\lambda^3\).  The identity
\[
 \frac{(1+\lambda Z)^3-1}{\lambda^3}
          =Z^3-\zeta^2\lambda Z^2-\zeta^2Z
\tag{A.170.11}
\]
reduces, when equated to one, to the separable irreducible polynomial
\(Z^3-Z-1\) over \(\F_3\).
Thus its root gives an unramified extension of degree three.
This representative has depth nine.

For later use, put \(b=1+\pi-\pi^2\).  The exact identity
\[
 \zeta-b^3=-\frac92-3\pi+\frac92\pi^3-3\pi^5
\tag{A.170.12}
\]
has valuation seven.  Consequently \(\zeta/b^3\) has depth seven.
The possible first depths of unit cubes are also explicit.
For \(v_{\mathfrak p}(w-1)=k\ge1\), the expansion of \(w^3-1\) gives
depth \(3k\) for \(k<3\), depth at least nine for \(k=3\),
and depth \(6+k\) for \(k>3\).
A cube with residue one has a cube root with residue one.
It follows that none of the depths one, five or seven can be
the first depth of a cube.  Multiplication by a power of the
depth-nine unramified representative cannot change these depths.
Since \(u_{11}u_{59}\) is a cube in \(F_{\mathfrak p}\), the local classes generated by
\(\epsilon,u_{11},u_{59},\zeta\), modulo the unramified line,
have the following exact kernel:
\[
 \epsilon^a u_{11}^b u_{59}^c\zeta^e
  \text{ is a cube times an unramified representative}
       \quad\Longleftrightarrow\quad
       a=e=0,\ b=c\quad\text{in }\F_3.
\tag{A.170.13}
\]
To see all cases directly, replace \(u_{59}\) by \(u_{11}^{-1}\)
modulo cubes and replace \(\zeta\) by \(\zeta/b^3\).
The first nonzero exponent then produces respectively depth
one, five or seven.  Each is excluded by the preceding cube
calculation.

\emph{The exact source of the anti-invariant classes.}
The roots of unity in \(F\) are \(\mu_6\): any cyclotomic subfield
is abelian over \(\Q\), and the abelianization of
\(\operatorname{Gal}(F/\Q)=S_3\) leaves only \(\Q(\zeta)\).
Dirichlet's theorem gives unit ranks two in \(F\) and one in
\(K_1\).  The norm-one unit group for \(F/K_1\) therefore has
free rank one and torsion \(\mu_6\).
The unit \(\epsilon\) is not a cube times a root of unity of
three-power order: its depth is one, whereas a power of \(\zeta\)
has depth at least three and a unit cube has depth at least three.
Its index in that free rank-one group is consequently prime to
three.  Hence \([\epsilon],[\zeta]\) are a basis of the
anti-invariant global unit classes modulo cubes.

For completeness, the passage to anti-invariant classes is literal.
The map on a class \([w]\) with
\(\tau[w]=-[w]\) which takes \(w\) to \((w/\tau w)^2\)
preserves that class modulo cubes and gives a representative of
norm one.  If a norm-one unit is a cube of a unit in \(F\), that
cube root has norm one: its norm is a cube root of unity in the
real field \(K_1\).

At each \(q=2,11,59\), the three primes of \(F\) are generated
by \(g_q,\sigma(g_q),\sigma^2(g_q)\); their norms are \(q^2\).
The involution fixes the first prime and interchanges the other
two.  Thus the \(u_q\) have the three independent anti-invariant
valuation vectors \((0,1,-1)\).  The three-prime is unique and
contributes no anti-invariant valuation.
Finally \(\operatorname{Cl}(F)\{3\}=0\), by A.138.
If an \(S_F\)-unit has all valuations divisible by three, dividing
its divisor by three gives a three-torsion ordinary ideal class,
which is zero.  Dividing by the resulting principal cube reduces
the class to a global unit.  This proves the complete basis
\[
          \mathcal A^-=
              \F_3[\zeta]\oplus\F_3[\epsilon]\oplus
              \F_3[u_2]\oplus\F_3[u_{11}]\oplus\F_3[u_{59}].
\tag{A.170.14}
\]
In this argument the anti-invariant valuation rows may first be
removed modulo three using the literal \(u_q\); the norm-one
projection just displayed makes the remaining unit class
anti-invariant.  No completeness assertion about a computed
unit group is used.

\emph{The root-stabilizer map at every height.}
Let \(L_n=\Q(\mu_{3^n},\rho_n)\).
Its global and local three-adic Galois groups are the full
holomorphs used in A.155, and its local extension at three is
totally ramified.  Above \(F\), the actual groups are
\[
 J=(3\Z/3^n)\rtimes U_1,\qquad
 U_1=\{v\in(\Z/3^n)^\times:v\equiv1\pmod3\},
 \qquad H=\operatorname{Gal}(L_n/E_n)=U_1.
\tag{A.170.15}
\]
Conjugating \(H\) by translation three produces the translations
\(3(1-v)\).  They generate \(9\Z/3^n\), since \(v=4\)
already gives its generator.  The normal closure of \(H\) in
\(J\) is therefore \((9\Z/3^n)\rtimes U_1\).
This proves that the largest Galois subfield of \(E_n/F\) is
exactly \(E_2=F(t)\), where \(t^3=d^2\alpha\).

The same calculation controls unramified pullback at three.
Let \(E_{n,\mathfrak P}\) be the completion at the unique three-prime.
The local splitting field is disjoint from the unramified cubic
extension \(F_{\mathfrak p}^{\rm ur,3}/F_{\mathfrak p}\).
Adjoining this field adds a direct \(C_3\) factor to A.170.15.
The same normal-closure calculation shows that every cyclic cubic
subfield over \(F_{\mathfrak p}\) inside
\(E_{n,\mathfrak P}F_{\mathfrak p}^{\rm ur,3}\) is in
\(E_{2,\mathfrak P}F_{\mathfrak p}^{\rm ur,3}\).
Since \(E_{n,\mathfrak P}/F_{\mathfrak p}\) is totally ramified,
its unique unramified cubic extension is
\(E_{n,\mathfrak P}F_{\mathfrak p}^{\rm ur,3}\).
Thus a cubic Kummer class over \(F_{\mathfrak p}\) becomes unramified
over \(E_{n,\mathfrak P}\) precisely when it is in the span of the
\(E_{2,\mathfrak P}/F_{\mathfrak p}\) radical and the unramified line.
The first radical has valuation two modulo three.
For a unit class its coefficient must be zero, leaving exactly
the local test in A.170.13.

At two the tower is unramified.  A nonzero coefficient of \(u_2\)
in A.170.14 retains a valuation prime to three at one of its two
support primes; the other four basis elements are units there.
Its Kummer extension is then ramified, so unramifiedness forces
that coefficient to zero.  At three A.170.13 and A.170.15 now
force precisely the line \(\F_3[u_{11}u_{59}]\).
Conversely the ramification checks in the next paragraph show
that line is everywhere unramified.  This proves A.170.6.

\emph{Nontriviality, ramification and the actual Frobenius values.}
At either changed prime the valuation vector of \(u\) in \(F\)
is \((0,1,-1)\).  The vector of \(d^2\alpha\) is \((2,2,2)\).
These do not span the same line modulo three.
By the largest-Galois-subfield calculation in A.170.15,
\(u\) is not a cube in any \(E_n\).  This proves the
nontriviality of A.170.2 and its degree three.
Its fixed field is generated by \(r+r^{-1}\) over \(K_n\),
which also proves \(\mathscr H_m=K_m\mathscr H_n\).

At eleven and fifty-nine the branches of \(E_n/F\) are totally
ramified of degree \(3^{n-1}\), for \(n\ge2\).
The valuations of \(u\) upstairs are therefore divisible by three.
At residue characteristic different from three, removing the
uniformizer cube leaves a unit Kummer extension, which is
unramified: its residue-field extension supplies the root of the
residue unit, and Hensel lifts it since three is a unit.
At other primes away three the radical is a unit.
At three it is already a cube in \(F_{\mathfrak p}\), by A.170.10.
Consequently \(E_n(r)/E_n\) is everywhere unramified.
Ramification in the cyclic cubic subfield over \(K_n\) is also
trivial, since its index divides three and, after this unramified
base change, divides \([E_n:K_n]=2\).

At every two-prime \(u\equiv1\pmod2\), by its last expression in
A.170.1.  The cube map is an automorphism of the principal units,
so these primes split in \(E_n(r)/E_n\), and hence in the
coprime-degree cubic subfield.  The same splitting holds at three.

At the changed linear branches, \(\alpha\equiv-2\pmod{11}\)
and \(\alpha\equiv-1/6\pmod{59}\).  Both substitutions give
\(u\equiv-\zeta\).  The residue calculations over \(E_n\) are
\[
 (-\zeta)^{(11^2-1)/3}=\zeta,\qquad
 (-\zeta)^{(59^2-1)/3}=\zeta^2 .
\tag{A.170.16}
\]
The residue degree of \(E_n/K_n\) at a linear branch is two.
Restriction of arithmetic Frobenius therefore squares its value
in the cyclic cubic group.  Inverting multiplication by two
in \(\Z/3\) gives the two \(K_n\)-values \(2,1\) in A.170.4.
This explicitly fixes the Frobenius convention and restriction map.

At level two the changed quadratic primes are principal by the
literal elements \(\sigma_{11},\sigma_{59}\) of A.165.11.
Their class characters are zero.  At every higher level, the
extension in A.170.2 is the same base change, so its character is
pulled back by the actual ideal norm.  The changed branches have
relative residue degree one, and retain value zero.
This finishes A.170.4 and compatibility of the characters.

The balanced relation at eleven is sent to
\(0-2\cdot2=2\) in \(\Z/3\), and the relation at fifty-nine
is sent to \(0-2\cdot1=1\).  Thus the precise character map does
not descend through \(C_n^{\rm bal}\).
The distinguished two-prime is in its kernel, but its ordinary
three-primary class is nonzero for \(n\ge3\) by A.168.
That kernel consequently has order at least three, while the
character is onto by A.170.4.  This proves A.170.5, also on the
subgroup generated by the \(S\)-prime classes.  This is a size
bound; it does not assert that the three-rank is at least two.

The exact field identities, norm-four unit construction, local
depths, finite residue characters and seven root-stabilizer
normal-closure calculations are retained in
\path{code/two_inert_genus_character_certificate.py}, with executed
output and data.  The independent degree-six reconstruction is in
\path{code/two_inert_anti_unit_independent_audit.py}, also with retained
output and data.  The proof above supplies the unramified and
class-group maps used with those finite calculations.
\end{proof}
